\section{APS}

\aps~is an analysis tool which generates high-level overview/summary pages of how a code performs.
The usage is close to trivial:

\begin{verbatim}
mpirun aps ./code args
#aps -report result folder
aps-report result folder
\end{verbatim}

\noindent
which should create a plain HTML file to study.
If you code does not use MPI, skip the \texttt{mpirun}.


The tool works best on Intel machinery, as it can read out the machine's performance counter.
These counters are not available (that way) on other vendors' machines.
In this case, you can restrict APS to focus on certain modi through \texttt{aps --collection-mode=mpi,openmp}.

\aps~provides a global overview of the code's behaviour.

Further to that, you might need \texttt{#SBATCH --disable-perfparanoid}


\paragraph*{Overview}
 
\begin{itemize}[noitemsep]
  \item[$\boxplus$] Simple to use with support for MPI, share memory parallelisation, I/O and vectorisation features.
  \item[$\boxplus$] Yields simple HTML pages which are easy to share and archive.
\end{itemize}
 
 
\begin{itemize}[noitemsep]
  \item[$\boxminus$] Provides the full range of information only on Intel CPUs.
  \item[$\boxminus$] Limited filter features to reduce the measurement overhead.
  \item[$\boxminus$] Cannot distinguish between different threads, and is not so
  straightforward to use with MPI. 
  \item[$\boxminus$] Focuses on timings and call counts only.
\end{itemize}

